% Part: history
% Chapter: biographies
% Section: georg-cantor

\documentclass[../../../include/open-logic-section]{subfiles}

\begin{document}

\olfileid{his}{bio}{can}

\olsection{গেয়র্গ কান্টর}

\olphoto{cantor-georg}{গেয়র্গ কান্টর}

গেয়র্গ কান্টরের (উচ্চারণ: গে-অর্গ কান-টোর) একটি
প্রথমদিকের জীবনীতে দাবি করা হয়েছিল যে রাশিয়ার সেন্ট
পিটার্সবার্গগামী জাহাজে তাঁর জন্ম হয় এবং সেখানেই তাঁকে
পাওয়া যায়; তাঁর বাবা-মায়ের পরিচয়ও নাকি অজানা ছিল।
দাবিটি অবশ্য সত্য নয়। তবে ১৮৪৫ সালে সেন্ট পিটার্সবার্গেই
তাঁর জন্ম হয়েছিল।

১৮৬৭ সালে কান্টর বার্লিন বিশ্ববিদ্যালয় থেকে গণিতে
ডক্টরেট লাভ করেন। সেটতত্ত্বে তাঁর কাজের জন্য তিনি পরিচিত;
স্বতন্ত্র গবেষণাক্ষেত্র হিসেবে সেটতত্ত্ব প্রতিষ্ঠার কৃতিত্বও
তাঁকে দেওয়া হয়। ভিন্ন আকারের অসীম সেট আছে, এটি
তিনিই প্রথম প্রমাণ করেন। তাঁর তত্ত্বগুলি, বিশেষ করে
অসীম নিয়ে তাঁর তত্ত্ব, সেই সময়ের গণিতবিদদের মধ্যে
বিস্তর বিতর্কের জন্ম দেয়; তাঁর কাজ ছিল বিতর্কিত।

কান্টরের ধর্মীয় বিশ্বাস ও গাণিতিক কাজ ওতপ্রোতভাবে
জড়িত ছিল। এমনকি তিনি দাবি করেছিলেন যে ট্রান্সফাইনাইট
সংখ্যার তত্ত্ব ঈশ্বর সরাসরি তাঁকে জানিয়েছেন। জীবনের শেষ
দিকে তিনি মানসিক অসুস্থতায় ভুগেছিলেন। ১৮৯৪ সাল থেকে
তাঁকে হাসপাতালে ভর্তি হতে হয়; পরবর্তী বছরগুলিতে তা
আরও ঘন ঘন ঘটে। গণিতবিদ লেওপোল্ড ক্রোনেকারের সঙ্গে
সম্পর্কচ্ছেদসহ তাঁর কাজের তীব্র সমালোচনা তাঁকে বিষণ্ণতা
ও গণিতে আগ্রহহীনতার দিকে ঠেলে দেয়। বিষণ্ণতার সময়ে
তিনি দর্শন ও সাহিত্যে মন দেন। শেক্সপিয়রের নাটকগুলির
প্রকৃত লেখক ফ্রান্সিস বেকন—এমন একটি তত্ত্বও তিনি
প্রকাশ করেছিলেন।

১৯১৮ সালের ৬ জানুয়ারি হালে-র একটি স্যানাটোরিয়ামে
কান্টরের মৃত্যু হয়।

\begin{reading}
কান্টরের পূর্ণাঙ্গ জীবনী জানতে \citet{Dauben1990} এবং
\citet{Grattan-Guinness1971} দেখো। তাঁর আমূল নতুন
মতগুলি বিবিসি রেডিও ৪-এর \emph{A Brief History of
  Mathematics} অনুষ্ঠানেও আলোচিত হয়েছে
\citep{Sautoy2014}। কান্টরের তত্ত্বগুলি র‍্যাপ সংগীতে
শুনতে চাইলে \citet{Rose2012} দেখো।
\end{reading}

\end{document}
